\documentclass[12pt, hyperref, a4paper, UTF8]{ctexart}

\usepackage{UCASReport}
\usepackage{enumitem}
\usepackage{booktabs}
\usepackage{datetime}

\setlist[1]{
    itemsep = 0pt,
    partopsep = 0pt,
    parsep = 0pt,
    topsep = 5pt,
    leftmargin = 2em,
    labelsep = 0.25em
}

\newcommand{\barline}[1]{\overline{#1\rule{0pt}{1.5ex}}}



\begin{document}

\cover
\thispagestyle{empty}
\clearpage
\setcounter{page}{1}



\section{\textit{On Being a Scientist} 的出版背景}
\textit{On Being a Scientist} 第三版是美国国家科学院、国家工程院和医学研究所联合出版，它的初衷是概述科学界的专业规范，此外重点突出现代科学界面临的新挑战，作者从科学精神、人类社会公众影响、科学家期待等三个角度出发，详细阐述了成为一名负责任的科学家应具备的素养和能力。这个作者从责任、诚信、态度、伦理等方面，给我们呈现了科学家应有的行为规范，并且对于学术不端行为进行全面剖析，给出解决策略。

目录大致分为 5 个部分，第一部分主要介绍负责任的科研行为基础与指导关系建立，第二部分聚焦于科学研究的具体实践与学术诚信，第三部分涵盖了学术违规行为的响应、受试者保护与实验室安全，第四部分探讨了研究成果公开、署名权与知识产权，第五部分延伸至科研人员在社会责任与价值观冲突，除上述核心规范外，\textit{On Being a Scientist} 最后还包含了一个案例研究讨论部分。

现如今，科学发展节奏极快且空前复杂，科研工作者往往缺乏足够时间与学生探讨专业标准，加之初学者也不一定总能得到最佳指导或模范行为引领，此外研究活动的跨学科和跨国家属性不断增强，学术界、工业界和政府社会之间联系日益紧密，而且数字化通信技术飞速发展，都引发许多全新的伦理实践挑战，\textit{On Being a Scientist} 正是在此背景下出版发布的，旨意为所有阶段的科研工作者提供负责任的科研行为指南，适合感兴趣的人阅读，也适合作为科普性读物，让更多人了解科学。



\section{学术不端的危害}
学术不端行为不仅是对科学核心价值观的严重违背，还会给个体和科学界带来毁灭性的后果，这种亵渎欺骗行为更多地表现在不同层面，造成极其恶劣且深远的危害，导致学术道德和科学文化的滑坡：

\begin{itemize}
    \item 第一条，对科学共同体和科研进展的破坏。科学是一项累积性事业，新的研究通常建立在过往的结果基础之上，如果研究者发表虚假数据，其他同行尝试复制或者延伸，将白白浪费大量宝贵时间和科研资金，这种不负责行为会阻碍研究领域进展，误入歧途，且之不端的行为会直接破坏研究者自身权威性和个人声音，危及职业生涯，最后由于科学结果深远地影响社会，虚假的结果会给公众带来直接危害，动摇社会对科学事业的信任。
    \item 第二条，对研究者及机构声誉的毁灭性打击。在科学界内部，这种行为带来的后果往往是毁灭性的，它会导致个人名誉扫地，学术生涯终止，还会让共事的同僚产生被深深背叛的感觉，而且一旦被实证存在篡改伪造剽窃，会连其所在的机构、共同作者以及整个研究领域的声誉遭受重创。
    \item 第三条，对公众利益的损害和社会信任的动摇。科学研究结果深刻地影响着社会，许多研究直接关系到个人的健康与福祉，比如医学治疗的危险性被虚假的评估，错误的信息进入科学记录，可能会对公众造成直接的身体伤害，此外，信息数字化发达的今天，误导性数据传播的速度和范围难以控制，其造成的后果更加不可预测，学术不端事件引发媒体和公众的负面关注，削弱科学界自我监督能力，动摇全社会对科学事业及其诚信的信任，最终甚至可能导致某些重要的科研领域失去原有的资金支持和公众支持。
\end{itemize}



\section{学术不端行为的表现形式}
学术不端行为通常指在提出、执行、审查研究或报告研究结果时出现的捏造、篡改、剽窃，统称为 FFP，具体包括有意和无意两种，有意的主要指未经授权或伪造数据、故意、重复发表、伪造同行评议、未经同行评议地发布论文，无意的主要指剽窃他人研究成果、引用不准确或不完整的文献资料、篡改数据图表文字以及不正当署名。学术不端行为还包括虚假陈述、不当署名等。

除了 FFP 之外，还存在许多有问题的研究实践，例如同行评审中违反保密原则、未合理分配论文署名权、隐瞒不报他人的学术不端行为、拆分发表等。不过，我们可以注意的是，构成学术不端的一个关键特征是其主观欺骗意图，这与科研中诚实的失误或由于疏忽造成的错误有着本质区别。但是学术不端行为破坏了科学研究等严肃性和可信度，而且对科研人员和科研机构造成不良的影响，甚至造成恶劣社会影响，为了避免学术不端行为的发生，需要国家和社会支持采取一些措施制止：

\begin{itemize}
    \item 科研工作者需要正确的道德教育；
    \item 需要对科研工作者进行学术规范培训；
    \item 对科研机构进行监督和管理。
\end{itemize}



\section{科研人员如何避免犯错误}
那么为了避免在科学研究中犯错，科研人员应采取以下有效措施：

\begin{itemize}
    \item 建立良好的指导体系，初级研究人员应积极寻求经验丰富的导师作为领路人，与导师确立清晰的沟通期望，从而在面临职业或道德抉择时获得正确的建议。
    \item 严谨客观地处理好数据，科研人员有责任义务极其详细准确地记录存储自己的工作流程，建立永久性的数据记录，确保其细节足以让其他人检查和复制实验。
    \item 保持开放的心态，承认错误，包容不同的思想，有勇气承认自己的错误，对科研诚实认真对待。
    \item 遵循科学的验证方法，科研需要承担风险且容易出错，但可以通过采用双盲试验、统计学检验、完善同行评审等科学手段来最小化错误率，及时纠正已存在的错误，如果发现由于客观限制或者人为疏忽而发表了错误的数据或结论，科研人员有责任通过在期刊上发表更正或勘误声明，来修正原始记录，防止误导后续科研工作者。
    \item 坚守诚信原则，科研工作者发表论文需遵守诚实守信的原则，不能伪造数据，保持对待科研热情的兴趣，将其作为事业融入自己的生活。
\end{itemize}



\section{科学研究中的道德准则和伦理规范}
作者在 \textit{On Being a Scientist} 一书提及，科学研究是一个过程，建立在诚实、公平、客观、公开、值得信赖和尊重他人的日常价值观基础之上，科研工作者履行着三大基本义务：尊重同行给予的信任、对自身的职业操守负责、服务公众的方式行事。

要考虑研究的目的、方式和范围，其次是研究对象，试验目的，实验方法，使用策略，对研究结果的影响，以及实验对社会国家的影响，研究对象可能会受到哪些影响，研究结果是否会引发社会不平等引起社会舆论，如果涉及动物实验，需要减少实验动物数量、优化程序过程、减轻动物痛苦、使用非活体材料替代活体动物，接受机构动物保护与使用委员会的审查。

作者认为科研工作者在学术上都是平等的，但不是所有人都有同等权利的资格。我们需要尊重学术自由原则、尊重同行评价原则、尊重知识产权原则、尊重个人隐私原则，还有作者提及科学研究成果是否公开、发表是否需要费用等一系列探讨。



\section{研究成果公开与知识产权}
在同行评审期刊上发表论文是传播公开研究成果最重要的方式，公开的成果要求研究者在借鉴他人想法时，必须通过极其谨慎的引文操作给予原作者适当的荣誉，绝对不能未经署名直接挪用，写论文之时，一般都是按照作者先后次序来决定，署名不仅是一种荣誉，也是一种责任，如果发现有什么问题，第一作者要对自己的作品内容负责，而且啊，署名一般要求我们那些在论文中做出了直接内容上贡献的科研工作者需被列入作者名单，尽可能对署名进行协商，从而最大限度地避免由于署名问题产生的矛盾，坚决抵制挂名荣誉作者或隐匿真实贡献幽灵作者等现象，商业机密通过保密来维持商业价值，研究人员有责任义务熟悉所在机构和资助方的知识产权分配政策，并尽早申请以避免影响学术成果的发表。其次，加强对知识产权的关注，既要对自己的技术成果进行保护，又要尊重他人的劳动成果，对发明创造来说，用专利权武装自己的知识产权。



\section{对未来的展望}
科学在本质上是一个高度依赖自我监督的共同体，面对日益复杂的信息化科研环境，科学界需要持续更新探讨科研伦理规范，科研工作者不仅要对自己和科学界责任，还要肩负着反思其研究成果将如何应用于更广泛社会的责任。

\textit{On Being a Scientist} 提出新的研究视角，将科学家的研究伦理视为科学素养，明确科研需要考虑的问题，让科研工作者意识到自己应该要做什么，未来咱们科研工作者继续在公共事务中发挥关键作用，既作为具有专业知识的顾问，也作为推动社会进步的普通公民，参与政策制定与公众教育，只要科学界继续坚守诚实、公平、开放和严谨等核心价值观，科学事业以及所服务的整个人类社会必须将持续保持繁荣与创造力。



\end{document}